% ============================================================
\chapter{用語集}\label{app:glossary}

\begin{longtable}{p{3.6cm}p{10.0cm}}
\toprule
用語 & 意味 \\
\midrule
\endfirsthead
\toprule
用語 & 意味 \\
\midrule
\endhead
\bottomrule
\endfoot
均衡リターン（$\Pi$） & 時価総額の比率が平均分散最適化の答えになるように逆算した期待リターン。$\Pi=\delta\Sigma\wmkt$。黙示のベータによる CAPM の期待リターンと同じ \\
逆最適化 & 比率から、その比率を最適にする期待リターンを求めること \\
リスク回避度（$\delta$、$\lambda$） & リスクをどれだけ嫌うかの係数。市場のリスクプレミアム ÷ 市場の分散。期間によらない \\
ビュー & 投資家の見方。期待リターン（またはリターン）の線形の組み合わせ $p_k\mu$ についての見込み $Q_k$ \\
絶対ビュー & 1 つの資産の期待リターンを言うビュー \\
相対ビュー & 資産の差を言うビュー。$p_k$ の和が 0 \\
ピック行列（$P$） & ビューの組み合わせを並べた行列 \\
ビューの不確かさ（$\Omega$） & ビューの外れの共分散。多くは対角 \\
$\tau$ & 均衡の不確かさの大きさ。事前分布 $\mu\sim\mathcal N(\Pi,\tau\Sigma)$ \\
採用率（$s$） & 均衡からビューへ動く割合。1 資産・1 ビューなら $\tau\sigma^2/(\tau\sigma^2+\omega)$ \\
事後の期待リターン（$\muBL$） & 均衡とビューを混ぜた期待リターン \\
情報形・更新形 & BL の公式の 2 つの書き方。精度を足す形と、均衡を驚きで更新する形 \\
市場版の BL & ビューを、期待リターンではなく来期に実現するリターンについてのものとみる版（Meucci）。$\tau$ が出てこない \\
He--Litterman の $\Omega$ & $\omega_k=\tau p_k\Sigma p_k\T$。ビュー 1 つなら採用率 50\% \\
Idzorek の確信度 & 0〜100\% の確信度から $\Omega$ を決める方法。閉じた式 $\omega_k=\tau p_k\Sigma p_k\T(1-C_k)/C_k$ \\
黙示の確信度 & 100\% 信じたときの比率の傾きに対して、実際に何割傾いたか \\
Meucci の確信度（$c$） & $\Omega=\frac1cP\Sigma P\T$ の $c$。エントロピー・プーリングでは事前と事後の確率を混ぜる割合 \\
情報係数（IC） & 予想と実際のリターンの相関。順位相関で測ることも多い \\
較正の傾き（$b$） & 実際を予想に回帰した傾き。予想の大きさまで含めた当たり具合。BL の採用率の目安 \\
Grinold の式 & 予想の点数を期待リターンに直す式。$\alpha=\IC\cdot\sigma\cdot z$ \\
基本法則 & 情報レシオ $\approx$ IC $\times\sqrt{\text{独立な予想の数}}$（Grinold and Kahn） \\
確信の積み上がり & 外れに共通の部分があるのに $\Omega$ を対角にすると、全体の水準への確信がビューの数に比例して強くなること \\
縮小推定 & 推定値を単純な目標に寄せて、推定の誤差を減らす方法。Ledoit--Wolf の共分散、Stein の平均など \\
ファクターモデル & リターンを少数のファクターと個別の部分で表すモデル。$r=Bf+\epsilon$ \\
露出（$B$） & 資産がファクターにどれだけ反応するか \\
EGARCH-M & 分散の対数を時系列でモデル化し（下落への非対称な反応を表せる）、平均の式に条件付きの標準偏差を入れたモデル \\
条件付き分散 & それまでの情報のもとでの、翌期のリターンの分散の予想 \\
GPR & 地政学リスク指数（Caldara and Iacoviello）。新聞の記事の割合から作る \\
EPU & 経済政策不確実性指数（Baker, Bloom and Davis）。新聞の記事の頻度から作る \\
モメンタム & 過去 12 か月（直近 1 か月を除く）によく上がった資産がしばらく上がりやすい傾向 \\
短期の反転 & 直近 1 か月によく上がった個別株が翌月に戻りやすい傾向 \\
モメンタムの急落 & 下げ相場のあとの反発で、モメンタムが大きく負けること \\
エントロピー・プーリング & シナリオの確率を、ビューを満たす中で事前にいちばん近く（相対エントロピーが最小に）なるように変える方法（Meucci） \\
相対エントロピー & 2 つの確率分布の違いの尺度。カルバック・ライブラーの情報量 \\
コピュラ & 多変量の分布のうち、各変数の分布を除いた「つながり方」の部分 \\
裾の依存 & 極端に大きく動くときの連動の強さ。正規分布のコピュラでは 0 に近づく \\
COP & コピュラ・オピニオン・プーリング（Meucci）。ビューの変数の分布を混ぜ、コピュラは保つ \\
VaR、CVaR & ある確率でしか超えない損失の水準と、それを超えたときの損失の平均 \\
整合的なリスク尺度 & 分散投資の効果を正しく表すための条件（劣加法性など）を満たすリスクの尺度 \\
トラッキングエラー（TE） & ベンチマークに対する超過リターンの標準偏差 \\
情報レシオ（IR） & 超過リターン ÷ TE \\
能動的な比率 & ポートフォリオの比率 − ベンチマークの比率 \\
実質銘柄数 & $1/\sum_iw_i^2$。比率が均等な何銘柄に当たるか \\
回転率 & 比率の入れ替わりの大きさ。$\frac12\sum_i|w_{i,t}-w_{i,t-1}|$ \\
融合 & 同じ資産への複数の予想と不確かさを 1 つにまとめること。精度加重、共分散の交差など \\
\end{longtable}
